% ============================================================
% section5_interactivity.tex - Раздел 5: Интерактивный выбор угла и работа со всеми углами
% ============================================================
% !TeX program = xelatex
% Компиляция: xelatex --shell-escape main.tex

\section{Интерактивная визуализация и обработка всех углов}

В предыдущих разделах мы научились загружать, усреднять и фильтровать терагерцовые сигналы, а также вычислять их спектры и пропускания для отдельного выбранного угла. Но в реальном эксперименте мы имеем дело с целым набором углов (например, от $0^\circ$ до $90^\circ$ с шагом в $15^\circ$). Анализировать каждый угол по отдельности, постоянно меняя код вручную, крайне неэффективно. 

В этом разделе мы решим эту проблему двумя профессиональными способами:
\begin{enumerate}
    \item Создадим удобный интерактивный интерфейс прямо в Matplotlib с помощью переключателей (RadioButtons), позволяющий на лету выбирать нужный угол и мгновенно перестраивать графики временного сигнала и спектра.
    \item Напишем функцию для автоматической пакетной обработки всех углов, интерполируем их на общую сетку частот и построим красивое семейство спектров пропускания на одном графике.
\end{enumerate}

\subsection{Создание элемента выбора угла}

Matplotlib — это не просто библиотека для статичной графики. В модуле \texttt{matplotlib.widgets} содержится набор простых элементов управления (кнопки, слайдеры, чекбоксы, переключатели), с помощью которых можно легко превратить обычный график в интерактивное микро-приложение.

Для выбора угла исследования идеально подходит класс \texttt{RadioButtons} (радиокнопки). Этот виджет позволяет выбрать ровно один вариант из предложенного списка.

\subsubsection{Использование RadioButtons из matplotlib.widgets}

Виджету радиокнопок для отрисовки требуется отдельная область рисования (ось, или \texttt{Axes}). Обычно для этого на основном поле графика выделяют небольшой прямоугольник сбоку.

Рассмотрим, как инициализировать этот виджет:

\begin{minipage}{\linewidth}
	\captionof{listing}{Инициализация виджета RadioButtons}
	\begin{minted}{python}
import matplotlib.pyplot as plt
from matplotlib.widgets import RadioButtons

# Создаем фигуру и ось для графика
fig, ax = plt.subplots(figsize=(10, 6))

# Выделяем область для кнопок [left, bottom, width, height] в координатах фигуры
ax_buttons = plt.axes([0.02, 0.4, 0.15, 0.25], facecolor='lightgray')

# Задаем список доступных углов (меток)
angles_list = ["0-90", "15-75", "30-60", "45-45", "60-30", "75-15", "90-0"]

# Создаем виджет RadioButtons
radio = RadioButtons(ax_buttons, angles_list)
	\end{minted}
\end{minipage}

Здесь метод \texttt{plt.axes([left, bottom, width, height])} создает новые независимые оси на той же фигуре. Координаты задаются в долях от размера окна (от 0 до 1). Параметр \texttt{facecolor} определяет цвет фона панели с кнопками.

\subsection{Обновление графиков при смене угла}

Для того чтобы график реагировал на клики пользователя, нам необходимо связать виджет с функцией обратного вызова (callback function).

\subsubsection{Реализация функции обратного вызова (Callback)}

При нажатии на кнопку Matplotlib автоматически вызывает зарегистрированную функцию и передает ей в качестве аргумента строковую метку выбранного элемента (например, \texttt{"15-75"}).

\warning{Многие начинающие программисты внутри функции обратного вызова полностью очищают оси с помощью \texttt{ax.clear()} и строят графики заново через \texttt{ax.plot()}. Это плохой тон! Полная очистка осей сбрасывает настройки сетки, подписи осей, границы графиков и вызывает сильное мерцание экрана. Намного профессиональнее и быстрее — сохранять ссылки на объекты линий и обновлять их данные с помощью метода \texttt{line.set\_ydata()}.}

\note{Обрати внимание на важную деталь, связанную с типами данных. Наш загрузчик \texttt{data\_loader} возвращает словарь \texttt{signals}, ключами которого являются кортежи чисел (например, \texttt{(10.0, 0.0)}). Однако виджет \texttt{RadioButtons} отображает элементы как строки и при клике передает в функцию обратного вызова именно строковое представление (например, строку \texttt{"(10.0, 0.0)"}). Если попытаться получить данные напрямую через \texttt{signals[label]}, Python выбросит исключение \texttt{KeyError}. Чтобы решить эту проблему, мы создаем вспомогательный словарь-преобразователь \texttt{label\_to\_key = \{str(k): k for k in sorted\_keys\}}, который сопоставляет строковую метку с исходным кортежем-ключом.}

Давай реализуем этот высокопроизводительный подход. Для этого нам потребуется создать новый Python-файл в корне нашего проекта. Создай в PyCharm файл с именем \texttt{interactive\_plots.py} рядом с остальными нашими модулями (\texttt{config.py}, \texttt{spectrum.py} и т.д.) и напиши в нём интерактивный визуализатор временных сигналов:

\needspace{4\baselineskip}
\captionof{listing}{Интерактивный визуализатор сигналов (interactive\_plots.py)}
\begin{minted}{python}
import numpy as np
import matplotlib.pyplot as plt
from matplotlib.widgets import RadioButtons
import config
import data_loader
import spectrum

def main_interactive():
    # 1. Загружаем экспериментальный датасет
    signals, backgrounds = data_loader.load_and_average_dataset(
        config.DATA_DIR
    )
    if not signals:
        print("Ошибка: данные для визуализации не найдены!")
        return

    # Сортируем оригинальные кортежи-ключи
    sorted_keys = sorted(list(signals.keys()))
    
    # Решение проблемы несоответствия типов: создаем отображение из строковых меток в исходные кортежи-ключи.
    # RadioButtons в Matplotlib отображает строковые представления элементов и передает строку в callback.
    label_to_key = {str(k): k for k in sorted_keys}
    labels = list(label_to_key.keys())
    
    initial_label = labels[0]
    initial_key = label_to_key[initial_label]

    # 2. Подготавливаем фигуру и оси
    # Оставляем место слева для панели управления
    fig, ax = plt.subplots(figsize=(11, 6))
    plt.subplots_adjust(left=0.25) # сдвигаем основной график вправо

    # Получаем данные для начального угла
    t_sig, E_sig = signals[initial_key]
    E_sig_dc = spectrum.remove_dc(E_sig)
    E_sig_win, win = spectrum.apply_gaussian_window(t_sig, E_sig_dc)

    # 3. Строим исходные линии и сохраняем ссылки на них
    line_raw, = ax.plot(
        t_sig, E_sig_dc, color='gray', alpha=0.5, 
        label='Исходный сигнал (без DC)'
    )
    line_win, = ax.plot(
        t_sig, E_sig_win, color='blue', linewidth=1.5, 
        label='Сигнал после окна'
    )
    
    ax.set_xlabel('Время задержки (пс)')
    ax.set_ylabel('Амплитуда (у.е.)')
    ax.set_title(
        f'Интерактивный анализ сигнала. Угол: {initial_label}'
    )
    ax.grid(True)
    ax.legend(loc='upper right')

    # 4. Создаем панель радиокнопок
    ax_radio = plt.axes(
        [0.02, 0.35, 0.16, 0.35], 
        facecolor='#f0f0f0'
    )
    radio = RadioButtons(ax_radio, labels)

    # 5. Функция обратного вызова для обновления данных
    def update_plot(label):
        # Получаем данные выбранного угла по его строковому представлению
        key = label_to_key[label]
        t_new, E_new = signals[key]
        E_new_dc = spectrum.remove_dc(E_new)
        E_new_win, _ = spectrum.apply_gaussian_window(
            t_new, E_new_dc
        )

        # Быстро обновляем данные по оси Y существующих линий
        line_raw.set_ydata(E_new_dc)
        line_win.set_ydata(E_new_win)

        # Обновляем заголовок графика
        ax.set_title(
            f'Интерактивный анализ сигнала. Угол: {label}'
        )
        
        # Пересчитываем лимиты по оси Y под новые данные
        ax.relim()
        ax.autoscale_view(scalex=False, scaley=True) # масштабируем только Y

        # Запрашиваем перерисовку холста
        fig.canvas.draw_idle()

    # Связываем переключатель с функцией обновления
    radio.on_clicked(update_plot)
    
    plt.show()

if __name__ == '__main__':
    main_interactive()
\end{minted}

Обрати внимание на синтаксис \texttt{line\_raw, = ax.plot(...)}. Запятая после имени переменной крайне важна — метод \texttt{ax.plot()} возвращает список объектов линий (даже если строится всего одна линия), а с помощью запятой мы распаковываем этот список, получая сам объект линии типа \texttt{Line2D}.

Вызов метода \texttt{ax.relim()} сообщает графику о необходимости пересчитать новые границы данных, а \texttt{ax.autoscale\_view()} плавно масштабирует видимую область по вертикали, чтобы новый импульс помещался целиком и не «вылезал» за границы кадра.

Интерфейс созданного нами интерактивного инструмента представлен на \figref{fig:5_1}.

\smartfigure{fig5.1.png}{Интерфейс с виджетом RadioButtons для интерактивного переключения исследуемых углов ТГц сигнала.}{fig:5_1}

\begin{minitaskbox}
    \textbf{Поэкспериментируй с виджетом!} \\
    Запусти файл \texttt{interactive\_plots.py}. Покликай по разным кнопкам углов. Обрати внимание, как плавно меняется форма импульса и как Гауссово окно автоматически находит и отслеживает центр максимального пика для каждого нового угла. 
    
    \textit{Задание на самостоятельную доработку:} Добавь на график отображение огибающей самого окна пунктирной красной линией (\texttt{'r--'}) и научи ее также обновляться при переключении углов.
\end{minitaskbox}

\subsection{Вычисление спектров для всех углов}

Интерактивный выбор удобен для визуального экспресс-анализа, но нам также необходимо уметь строить интегральные картины зависимости пропускания от угла поворота поляризатора. Для этого нам нужно пройтись циклом по всем загруженным углам, рассчитать для каждого из них спектр пропускания и собрать результаты в единую структуру данных.

\subsubsection{Реализация функции compute\_all\_spectra}

Здесь возникает важный нюанс: теоретически, временные сетки для разных углов могут немного различаться (например, иметь разную длину или мелкие сдвиги из-за погрешностей позиционирования каретки задержки в спектрометре). Чтобы все спектры можно было корректно сопоставлять и хранить в единой двумерной матрице, их нужно интерполировать на единую сетку частот.

Давай реализуем функцию пакетного вычисления спектров в файле \texttt{spectrum.py}:

\needspace{4\baselineskip}
\captionof{listing}{Пакетная обработка и вычисление всех спектров (spectrum.py)}
\begin{minted}{python}
def compute_all_spectra(signals, backgrounds, pad_factor=config.PAD_FACTOR):
    """
    Вычисляет спектры пропускания для всех доступных углов.
    Возвращает словарь вида {угол: (freqs, transmission)}.
    """
    results = {}
    
    # Сортируем углы для упорядоченной обработки
    sorted_angles = sorted(list(signals.keys()))
    
    for angle in sorted_angles:
        t_sig, E_sig = signals[angle]
        
        # Находим соответствующий фон в словаре backgrounds
        # Если точного совпадения нет, берем первый попавшийся фон
        if angle in backgrounds:
            t_bg, E_bg = backgrounds[angle]
        else:
            first_bg_key = list(backgrounds.keys())[0]
            t_bg, E_bg = backgrounds[first_bg_key]
            
        # Рассчитываем пропускание с помощью нашего конвейера
        freqs, _, _, transmission = calculate_transmission(
            t_sig, E_sig, t_bg, E_bg, pad_factor
        )
        
        results[angle[0]] = (freqs, transmission)
        
    return results
\end{minted}

Мы возвращаем словарь, в котором ключами являются углы поворота первого поляризатора (например, \texttt{15.0}), а значениями — кортежи с индивидуальной шкалой частот и массивом коэффициента пропускания.

\subsection{Визуализация семейства спектров}

Теперь, когда у нас есть рассчитанные спектры пропускания для всех углов, мы можем визуализировать их все вместе на одном графике. Это позволит нам наглядно увидеть, как именно вращение поляризаторов меняет спектральные характеристики проходящего терагерцового излучения.

Напишем код для отображения семейства спектров в конце нашего скрипта \texttt{interactive\_plots.py}:

\needspace{4\baselineskip}
\captionof{listing}{Визуализация семейства спектров пропускания (interactive\_plots.py)}
\begin{minted}{python}
def plot_all_transmissions():
    # Загружаем данные и вычисляем спектры для всех углов
    signals, backgrounds = data_loader.load_and_average_dataset(config.DATA_DIR)
    all_spectra = spectrum.compute_all_spectra(signals, backgrounds)
    
    plt.figure(figsize=(10, 6))
    
    # Используем цветовую карту для красивого плавного градиента линий
    colors = plt.cm.plasma(np.linspace(0, 0.85, len(all_spectra)))
    
    for (angle, (freqs, transmission)), color in zip(all_spectra.items(), colors):
        # Строим график только в полезном диапазоне частот
        mask = (freqs >= config.F_MIN) & (freqs <= config.F_MAX)
        plt.plot(freqs[mask], transmission[mask], label=f'Угол {angle}', color=color, linewidth=1.5)
        
    plt.xlim(config.F_MIN, config.F_MAX)
    plt.ylim(0, 1.1)
    plt.xlabel('Частота (ТГц)')
    plt.ylabel('Коэффициент пропускания')
    plt.title('Семейство спектров пропускания для различных углов поворота поляризатора')
    plt.grid(True, linestyle='--', alpha=0.7)
    plt.legend(loc='lower left', ncol=2) # выводим легенду в две колонки
    plt.tight_layout()
    plt.show()
\end{minted}

В этом листинге мы использовали очень полезный трюк с цветовой картой (\texttt{colormap}): \texttt{plt.cm.plasma(np.linspace(0, 0.85, len(all\_spectra)))} генерирует массив цветов, плавно переходящих от сине-фиолетового к ярко-оранжевому. Это делает графики гораздо более эстетичными и легко читаемыми по сравнению со стандартным набором цветов Matplotlib.

Кроме того, мы применили логическую маску \texttt{mask} для фильтрации частот. Это позволило нам не только настроить границы отображения с помощью \texttt{plt.xlim()}, но и гарантировать, что Matplotlib не будет тратить ресурсы на отрисовку сотен невидимых точек вне рабочего частотного диапазона.

Результат построения семейства спектров показан на \figref{fig:5_2}.

\smartfigure{fig5.2.png}{Семейство спектров пропускания исследуемого образца при различных углах поворота поляризатора.}{fig:5_2}

\subsection{Полный код модуля}

Для закрепления материала и решения поставленной мини-задачи полный код файла \texttt{interactive\_plots.py} вынесен в \textbf{Приложение \ref{app:interactive_plots}}. В этот листинг интегрировано решение по динамическому отображению и автоматическому масштабированию огибающей Гауссова окна для каждого выбранного угла.

\subsection{Итоги раздела}

В этом разделе мы совершили качественный скачок от разрозненных файлов к целостному пакетному анализу данных:
\begin{enumerate}
    \item Овладели инструментами создания интерактивных интерфейсов в Matplotlib с помощью виджета \texttt{RadioButtons}.
    \item Реализовали производительный алгоритм обновления данных на графиках без мерцания и сброса осей через обновление объектов \texttt{Line2D}.
    \item Написали функцию \texttt{compute\_all\_spectra} для автоматического расчёта спектральных характеристик по всему датасету.
    \item Создали красивую интегральную визуализацию семейства спектров пропускания с градиентным окрашиванием линий.
\end{enumerate}

\note{\textbf{Как изменился проект в этом разделе:}
\begin{itemize}
    \item В модуле \texttt{spectrum.py} появилась новая функция \texttt{compute\_all\_spectra(signals, backgrounds, pad\_factor)}, автоматизирующая полный цикл расчётов для всех доступных в датасете углов.
    \item Был создан новый скрипт \texttt{interactive\_plots.py}, содержащий интерактивный графический инструмент \texttt{main\_interactive()} для инспекции импульсов и вспомогательную функцию \texttt{plot\_all\_transmissions()} для визуализации семейства кривых пропускания.
\end{itemize}
}

Теперь всё готово к тому, чтобы начать физический анализ результатов! В следующем разделе мы подключим мощный теоретический аппарат — модель угловой зависимости Бланко, и научимся сопоставлять теорию с полученными экспериментальными точками.
